\documentclass[12pt,a4paper]{article}
\usepackage{graphicx}
\usepackage{amsmath}
\usepackage{hyperref}
\usepackage{geometry}
\geometry{margin=2.5cm}

\title{\textbf{CDFD Runtime Paper 10: Multi-Domain Isomorphism and Adapter Evidence}}
\author{Steve Bico Mujjabi, MD\\
\small Independent Researcher\\
\small ORCID: \href{https://orcid.org/0009-0001-0556-5516}{0009-0001-0556-5516}}
\date{May 2026}

\begin{document}
\maketitle

\begin{abstract}
\noindent \textbf{Executive synthesis.} Multi-domain isomorphism is the claim
that many systems can be mapped into a common flow-constraint grammar. The
current runtime exposes 196 registered domain adapters through the CLI. This is
evidence of software breadth, not empirical proof that all domains obey CDFD.
This paper explains the adapter contract and connects it to the finance-climate,
universal-collapse, drug-synergy, and frontier-sweep experiments.
\end{abstract}

\begin{figure}[ht]
\centering
\fbox{\begin{minipage}{0.92\textwidth}
\centering
\textbf{Runtime evidence path}\\[0.4em]
Prior CDFD mathematics $\longrightarrow$ CDFL/runtime state $\longrightarrow$ bounded output $\longrightarrow$ falsification target.
\end{minipage}}
\caption{Conceptual schematic for the runtime papers. The engine translates the mathematical frame from the prior CDFD papers into executable CDFL/runtime diagnostics; candidate outputs remain tied to explicit tests before any stronger claim is made.}
\end{figure}

\section{Prior CDFD Basis}
This manuscript does not rederive the mathematical core of CDFD. It uses the
Mujjabi Adaptive Operating Ratio and the constrained-vacuum notation introduced
in Part I \cite{MujjabiI2026}, the torus-knot hierarchy and boundary-mode work
from the Part I sequence \cite{MujjabiVI2026,MujjabiIX2026}, the public balance
law developed in the Part I capstone \cite{MujjabiX2026}, the Part I transport
tests \cite{MujjabiXII2026}, and the Life Number and tri-regime biology bridge
from Part II \cite{MujjabiPartII2026}. The runtime papers therefore act as an
execution layer: they keep the mathematics connected to commands, outputs,
falsification tests, and domain-specific limits.


\section{The Adapter Contract}
Every domain adapter answers three questions:
\begin{enumerate}
    \item What is the domain expression of flow $\Phi$?
    \item What is the domain expression of constraint $C$?
    \item What interpretation is justified for $\Psi_s=(\Phi/C)S M_s$?
\end{enumerate}
If these questions cannot be answered with measurable proxies, the adapter is
not scientifically mature.

\section{Runtime Breadth}
The command
\begin{verbatim}
python cdfd.py domains
\end{verbatim}
currently lists 196 registered domain keys. Domains include medicine,
physics, origins of life, ecology, economics, climate, infrastructure,
materials science, software engineering, and many more.

This breadth is useful for hypothesis generation. It does not replace
domain-specific validation.

\section{Case Study: Finance and Climate}
\texttt{experiments/notebooks/discovery\_finance\_climate.py} models a coupled
toy system and tests whether financial $\Psi$ volatility precedes increased
climate constraint. The report presents a candidate directed edge. The key
control is to remove coupling. If the edge persists after decoupling, the
result is an artifact.

\section{Case Study: Universal Collapse}
\texttt{experiments/run\_universal\_discovery.py} models a central overloaded
node in a network. It records hub $\Phi$, hub $C$, hub $\Psi$, system mean
$\Psi_s$, and collapsed-node count in
\texttt{experiments/outputs/universal\_collapse.json}. This is a useful generic
cascade diagnostic for power grids, markets, organ systems, and transport
networks, but each real domain still requires its own data mapping.

\section{Case Study: Pharmacological Synergy}
\texttt{experiments/run\_drug\_synergy\_predictions.py} compares same-axis
interventions with cross-axis interventions. In CDFD terms, one drug lowering
$\Phi$ and another improving recovery or relaxation may behave differently from
two drugs both lowering $\Phi$. This is a toy pharmacology diagnostic, not a
prescribing rule.

\section{Frontier Sweep as Adapter Triage}
\texttt{experiments/outputs/frontier\_sweep.h5} records 2,500 parameter points
and 2,487 non-trivial candidate records. The proper use is triage: identify
which parameter regions deserve deeper testing, not declare thousands of laws.

\section{Failure Conditions}
Multi-domain isomorphism fails locally if:
\begin{itemize}
    \item the adapter cannot define measurable $\Phi$ and $C$;
    \item predicted thresholds do not reproduce;
    \item the model performs no better than domain baselines;
    \item the same parameter story is forced onto incompatible systems.
\end{itemize}

\section{Conclusion}
The adapter library is an engine of disciplined analogy. It lets one runtime ask
similar questions across many domains. The scientific standard remains local:
each adapter must earn its interpretation with data, reproducible diagnostics,
and falsification.
\begin{thebibliography}{9}
\bibitem{MujjabiI2026}
Steve Bico Mujjabi, MD. \emph{Vortex Stability in a Constrained Transport Vacuum and the Origin of the Fine-Structure Constant}. CDFD Part I Paper I, preprint, 2026.

\bibitem{MujjabiVI2026}
Steve Bico Mujjabi, MD. \emph{The Universal Torus Knot Hierarchy: $Z_n$ Symmetry, the Cinquefoil Family, and the General Structure of CDFT Particle Families}. CDFD Part I Paper VI, preprint, 2026.

\bibitem{MujjabiIX2026}
Steve Bico Mujjabi, MD. \emph{Even-$n$ Torus Knots in CDFT: The Exclusion of $T(2,2)$, the Extension to $T(2,2m)$, and the Anti-Phase Mass Scaling Law}. CDFD Part I Paper IX, preprint, 2026.

\bibitem{MujjabiX2026}
Steve Bico Mujjabi, MD. \emph{The Vacuum Equation of State: Deriving Torus Knot Self-Consistency from the Public CDFD Balance Law $\Psi = \Phi/C$}. CDFD Part I Paper X, preprint, 2026.

\bibitem{MujjabiXII2026}
Steve Bico Mujjabi, MD. \emph{Friction, Superconductivity, Plasma States, and Transport Mysteries: Observable Tests of Adaptive Constraint}. CDFD Part I Paper XII, preprint, 2026.

\bibitem{MujjabiPartII2026}
Steve Bico Mujjabi, MD. \emph{CDFD Part II: Origins of Life and Tri-Regime Bioenergetics}. Public release archive, 2026, \url{https://doi.org/10.5281/zenodo.20250821}.
\end{thebibliography}
\end{document}
